% Source: https://artofproblemsolving.com/community/c3h2419019_square_root_question
% Stable identity: aops-topic-2419019-q01
For $x\ne0$,
\[1-\dfrac{x-1}{x}=\dfrac{x-(x-1)}{x}=\dfrac1x,\qquad\text{so}\qquad\dfrac{x}{1-\dfrac{x-1}{x}}=\dfrac{x}{\dfrac1x}=x^2.\]
Hence the expression equals $\sqrt{x^2}$.

The symbol $\sqrt{\ \ }$ always means the \emph{nonnegative} square root, so $\sqrt{x^2}=|x|$ for every real $x$. (For example $\sqrt{(-3)^2}=\sqrt9=3$, not $-3$.) Since $x>0$, $|x|=x$, and the expression simplifies to $\boxed{x}$.

The sign $\pm$ appears when we \emph{solve} an equation such as $t^2=x^2$, which has the two solutions $t=\pm x$; but a single expression such as $\sqrt{x^2}$ has only one value. That is why the condition $x>0$ is given: it lets us replace $|x|$ by $x$.
